\documentclass[12pt]{article}
\usepackage[utf8]{inputenc}
\usepackage[T1]{fontenc}
\usepackage{hyperref}
\title{Wangkayutyuru Language Resources for K-12 Curriculum Development}
\author{Compiled by Claude Sonnet 5.5 (claude-sonnet-5-5) for Andr\'as Kornai}
\date{}
\begin{document}
\maketitle

\section{Introduction}
Wangkayutyuru (ISO 639-3 wky) is an extinct Australian language. The search results suggest that it is an alternate name for Wanggamala (also written Wangka-Yutjurru), traditionally spoken in the Northern Territory around the Hay River, east of Alice Springs. One source reports a single remaining speaker as of 2003. This identification comes from a search summary, and I did not confirm it against the pages themselves. It should not be confused with the living Western Desert varieties Wangkajunga and Wangkatja, which are separate. Very little exists for Wangkayutyuru.

\section{Primary Educational Resources}
I found no primers, textbooks or school programmes for Wangkayutyuru.

\section{Linguistic and Reference Materials}
\begin{itemize}
\item PHOIBLE sound inventory page carrying the Wangkayutyuru label, filed under Wanggamala: \url{https://phoible.org/languages/wang1289}
\item PHOIBLE source record (Jones 2011): \url{https://phoible.org/sources/319249}
\item Wikipedia entry for Wanggamala (ISO 639:lnw): \url{https://wikipedia.com/wiki/ISO_639:lnw}
\end{itemize}
I did not open these pages, so I cannot say what each contains beyond its title.

\section{Religious and Community Resources}
None found for Wangkayutyuru.

\section{Digital and Audio Resources}
None found for Wangkayutyuru.

\section{Curriculum Development Context}
A heritage curriculum is not possible from the materials I found. Any work would need archival research, for example at AIATSIS, and consultation with descendant communities in the Northern Territory. The nearby living Western Desert varieties are a different language and should not be substituted for it.
\end{document}